% GENERATED FILE. Edit canonical lessons, then run npm run generate:books.
% canonical-source-hash: fnv1a64:90c325350bb177f4
% canonical-lessons: FA-C18-yazdah, FA-C18-sizdah, FA-C18-panzdah, FA-C18-bist-sad, FA-C18-practice

\chapter{The Teens Are a Rule, Not Nine Words}
\label{ch:fa-teens}
\hlchaptermodality{18}

\begin{chapteropening}
\textbf{By the end of this chapter:} \emph{I can build any number from eleven to nineteen by putting a unit in front of dah, say twenty and a hundred, and write the tenth letter of this book.}

Three teens taught and six built from the rule, then the two round amounts a price actually names — and the one letter a hundred needs that no lesson yet teaches the hand to make.
\end{chapteropening}

\section[yâzdah]{\fa{یازده} — eleven, and one rule for nine numbers}

\label{lesson:FA-C18-yazdah}



\begin{quote}
Count one to ten, then say \textbf{dah} on its own and what it means.
\end{quote}

\subsection*{You'll want to know first — one word, and then a rule}
\begin{quote}
\textbf{\fa{یازده}} — \emph{yâzdah} — \textbf{eleven}
\end{quote}

Say it slowly and the ten is audible at the end: \emph{yâz} + \textbf{dah}. The first piece is \emph{yek} worn down; the second is the word for ten you learned in the last lesson.

\textbf{That is the whole of the teens.} Persian puts the unit in front of ten and stops:

\noindent
\begin{tabularx}{\linewidth}{@{}>{\raggedright\arraybackslash}X>{\raggedright\arraybackslash}X@{}}
\toprule
\textbf{} & \textbf{} \\
\midrule
12 & \emph{davâz}-\textbf{dah} — from \emph{do} \\
14 & \emph{chahâr}-\textbf{dah} — \emph{chahâr} unchanged \\
19 & \emph{nuz}-\textbf{dah} — from \emph{noh} \\
\bottomrule
\end{tabularx}

Nine numbers, one rule, and the unit is a number you already have. This is why ten had to come first.

\subsection*{Your turn}
\begin{itemize}
\raggedright
  \item \emph{Say it:} \textbf{yâzdah} — eleven
  \item \emph{Split:} \emph{yâz} + \emph{dah}
  \item \emph{Build:} \emph{chahâr} + \emph{dah} $\to$ \textbf{chahârdah}, fourteen
\end{itemize}

\begin{tcolorbox}[breakable,colback=teal!4,colframe=teal!35!black,title={Before you move on}]
Sources: \href{https://en.wiktionary.org/wiki/\%DB\%8C\%D8\%A7\%D8\%B2\%D8\%AF\%D9\%87}{Wiktionary: \fa{یازده}}; \href{https://www.omniglot.com/language/numbers/persian.htm}{Omniglot, Numbers in Persian}.
\end{tcolorbox}
\section[sizdah]{\fa{سیزده} — thirteen, and how far a unit can bend}

\label{lesson:FA-C18-sizdah}



\begin{quote}
Say \textbf{yâzdah}, split it in two, and build fourteen.
\end{quote}

\subsection*{You'll want to know first — one word}
\begin{quote}
\textbf{\fa{سیزده}} — \emph{sizdah} — \textbf{thirteen}
\end{quote}

\emph{si} + \textbf{dah}. The unit is \emph{se}, with its vowel lengthened and darkened — and English does the same thing without anyone noticing: \emph{three} becomes \emph{thir}- in \textbf{thirteen}, and \emph{five} becomes \emph{fif}- in \textbf{fifteen}.

Listen past the bending, then, and hear what does not bend: \textbf{the ten stays the same}, on the end, every time.

\subsection*{Your turn}
\begin{itemize}
\raggedright
  \item \emph{Say it:} \textbf{sizdah} — thirteen
  \item \emph{Split:} \emph{si} + \emph{dah}
  \item \emph{Compare:} \textbf{se} $\to$ \emph{si}, as English \textbf{three} $\to$ \emph{thir}-
\end{itemize}

\begin{tcolorbox}[breakable,colback=teal!4,colframe=teal!35!black,title={Before you move on}]
Sources: \href{https://en.wiktionary.org/wiki/\%D8\%B3\%DB\%8C\%D8\%B2\%D8\%AF\%D9\%87}{Wiktionary: \fa{سیزده}}; \href{https://www.omniglot.com/language/numbers/persian.htm}{Omniglot, Numbers in Persian}.
\end{tcolorbox}
\section[pânzdah]{\fa{پانزده} — fifteen, and the one irregularity worth naming}

\label{lesson:FA-C18-panzdah}



\begin{quote}
Say \textbf{sizdah} and split it, then say \textbf{panj}.
\end{quote}

\subsection*{You'll want to know first — one word}
\begin{quote}
\textbf{\fa{پانزده}} — \emph{pânzdah} — \textbf{fifteen}
\end{quote}

\emph{pân} + \emph{z} + \textbf{dah}. There is an \textbf{n} in there that \emph{panj} does not have in the place you would expect it, and there is a \textbf{z} that came from the \emph{j}.

Sixteen does the same thing: \emph{shânzdah}, from \emph{shesh}. Fifteen and sixteen are the only two of the nine teens that bend this far, and they bend the same way — so it is one irregularity, met once, not two.

\subsection*{Your turn}
\begin{itemize}
\raggedright
  \item \emph{Say it:} \textbf{pânzdah} — fifteen
  \item \emph{Say it:} \textbf{shânzdah} — sixteen, the same bend
  \item \emph{Run through:} \textbf{yâzdah, sizdah, pânzdah} and the ten inside each
\end{itemize}

\begin{tcolorbox}[breakable,colback=teal!4,colframe=teal!35!black,title={Before you move on}]
Sources: \href{https://en.wiktionary.org/wiki/\%D9\%BE\%D8\%A7\%D9\%86\%D8\%B2\%D8\%AF\%D9\%87}{Wiktionary: \fa{پانزده}}; \href{https://www.omniglot.com/language/numbers/persian.htm}{Omniglot, Numbers in Persian}.
\end{tcolorbox}
\section[bist sad]{\fa{بیست} and \fa{صد} — twenty and a hundred}

\label{lesson:FA-C18-bist-sad}



\begin{quote}
Say \textbf{pânzdah} and \textbf{sizdah}, and split each in two.
\end{quote}

\subsection*{You'll want to know first — two words}
\begin{quote}
\textbf{\fa{بیست}} — \emph{bist} — \textbf{twenty}
\end{quote}

\begin{quote}
\textbf{\fa{صد}} — \emph{sad} — \textbf{a hundred}
\end{quote}

Two words, taken together because a price names one or the other far more often than it names any of the eighteen numbers between them.

\textbf{\fa{صد}} is two letters and it is the largest number in this book. Its first letter, \textbf{\fa{ص}}, is one no lesson has taught the hand to make — it enters this book by eye only, the way most of the track's words have, and reading is what the number needs.

\begin{cousinweb}
\textbf{\fa{بیست}} goes back to Indo-European *\emph{wikmtih}, \textquotedblleft{}twenty\textquotedblright{} — Latin \textbf{vīgintī}, Sanskrit \emph{vimshati}.

\textbf{\fa{صد}} goes back to *\emph{kmtom}, \textquotedblleft{}hundred\textquotedblright{} — Latin \textbf{centum}, Greek \textbf{hekatón}, English \textbf{hundred}. The Latin form is the one English borrowed twice over, in \textbf{century} and \textbf{per cent}, so a reader who has ever seen a percentage has already met \textbf{sad}'s cousin.

The spelling is the joke: \textbf{\fa{صد}} is written with a borrowed Arabic letter for a word that is not Arabic at all.
\end{cousinweb}

\subsection*{Your turn}
\begin{itemize}
\raggedright
  \item \emph{Say it:} \textbf{bist} — twenty; \textbf{sad} — a hundred
  \item \emph{Connect:} \textbf{sad} $\leftarrow$ *\emph{kmtom} $\to$ Latin \textbf{centum}, English \textbf{hundred}
  \item \emph{Run through:} \textbf{dah, bist, sad} — the three round amounts
\end{itemize}

\begin{tcolorbox}[breakable,colback=teal!4,colframe=teal!35!black,title={Before you move on}]
Sources: \href{https://en.wiktionary.org/wiki/\%D8\%B5\%D8\%AF}{Wiktionary: \fa{صد}}; \href{https://www.omniglot.com/language/numbers/persian.htm}{Omniglot, Numbers in Persian}.
\end{tcolorbox}
\section[Practice]{Payoff — one to twenty, and the hundred above it}

\label{lesson:FA-C18-practice}



\begin{quote}
Count one to ten, then say a hundred.
\end{quote}

\subsection*{Your turn}
Three of the nine teens were taught. \textbf{Build the other six} rather than recalling them.

\begin{enumerate}
\raggedright
  \item From \emph{do}, \emph{chahâr}, \emph{haft}, \emph{hasht} and \emph{noh}, build twelve, fourteen, seventeen, eighteen and nineteen by putting the unit in front of \textbf{dah}. Then check each against a printed list and repair only what missed.
  \item Say sixteen, and say which taught number it bends like.
  \item Count one to twenty straight through.
  \item Hear twelve amounts named singly, from \textbf{yek} to \textbf{sad}, and say each.
\end{enumerate}

The point of the chapter is that six of these nine numbers were never taught and you can say them anyway. That is what a rule buys and a word list does not.

\begin{tcolorbox}[breakable,colback=teal!4,colframe=teal!35!black,title={Before you move on}]
Source: \href{https://www.omniglot.com/language/numbers/persian.htm}{Omniglot, Numbers in Persian}.
\end{tcolorbox}
